\subsection{محدودیت‌ها و تحلیل شکست}
با وجود کارایی قابل توجه سخت‌افزاری و موفقیت در \lr{DistilBERT}، نقاط ضعفی نیز در عملکرد سایر مدل‌ها مشاهده می‌شود:
\begin{enumerate}
    \item \textbf{انقباض تفکیک‌پذیری در ابعاد بزرگ‌تر (\lr{PASS\_LOW} در \lr{RoBERTa} و \lr{GPT-2}):} عملکرد ماژول در مدل \lr{RoBERTa} روی ۵۰.۹۲٪ و در \lr{GPT-2} روی ۵۵.۵۰٪ قرار گرفته است. تحلیل شکست نشان می‌دهد که با افزایش ابعاد سر ($d \ge 64$)، نگاشت اشباع‌شده تانژانت هذلولوی درایه‌ها را به مقادیر مرزی $\pm 1$ فشرده کرده و تنوع زاویه‌ای فضای تعبیه را کاهش می‌دهد؛ در نتیجه، در غیاب یک پیش‌شرط‌ساز طیفی، ظرفیت تفکیک توکن‌های نحوی در لایه‌های ابتدایی افت می‌کند.
    \item \textbf{افت دقت صفر-شات در \lr{Qwen} (\lr{48.74\%}):} در مدل‌های زبانی عظیم، تکیه بر نگاشت شبه‌باینری بدون بازآموزی هم‌زمان وزن‌های پروجکشن اصلی باعث تخریب کالیبراسیون توزیع لاجیت‌ها گردیده و دقت خروجی را در حاشیه شانس تصادفی قرار می‌دهد.
\end{enumerate}
